\documentclass[11pt]{article}

\usepackage[T1]{fontenc}
\usepackage[utf8]{inputenc}
\usepackage{lmodern}
\usepackage{amsmath,amssymb,amsthm,mathtools}
\usepackage[a4paper,margin=28mm]{geometry}
\usepackage{microtype}
\usepackage{xurl}
\usepackage[hidelinks]{hyperref}
\usepackage{enumitem}
\usepackage{booktabs}
\usepackage{array}

\newtheorem{theorem}{Theorem}[section]
\newtheorem{lemma}[theorem]{Lemma}
\newtheorem{proposition}[theorem]{Proposition}
\newtheorem{corollary}[theorem]{Corollary}
\newtheorem{definition}[theorem]{Definition}
\newtheorem{remark}[theorem]{Remark}

\newcommand{\N}{\mathbb N}
\newcommand{\WIN}{\ensuremath{\mathrm{WIN}}}
\newcommand{\LOSS}{\ensuremath{\mathrm{LOSS}}}
\newcommand{\DRAW}{\ensuremath{\mathrm{DRAW}}}
\newcommand{\odd}{\operatorname{odd}}
\newcommand{\src}{\operatorname{src}}
\newcommand{\ct}{\operatorname{ct}}
\newcommand{\asl}{\operatorname{asl}}
\newcommand{\Moves}{\operatorname{moves}}

\title{The Two-Player $3n\mathbin{\pm}1$ Game:\
A Binary Normal Form and a Well-Founded No-Draw Proof}
\author{Grisha Pochuev\\
  \small\href{mailto:n_854@mail.ru}{n\_854@mail.ru}\\
  \small\url{https://github.com/Grisha-Pochuev/3n-plus-minus-1-game}}
\date{10 August 2026}

\begin{document}
\maketitle

\begin{abstract}
Consider the impartial two-player game on positive odd integers in which a
move from $n>1$ chooses $3n-1$ or $3n+1$ and removes all powers of two; reaching
$1$ wins immediately. Arbitrary play need not terminate: for example
$5\to 7\to 5$ is a legal cycle. The question is whether optimal play nevertheless
has finite remoteness from every starting position. We conjugate the game to a
binary normal form with an expanding move $A$ and a strictly contracting move
$B$, introduce constant-tail coordinates and a three-free source map, and rank
marked proof configurations by retained source, finite proof-height tokens, and
a typed fixed-fibre normalizer. The key entry step is a factor-removal lemma
that either removes one factor of three at the same tail exponent or produces
finite outcome provenance. This prevents an unranked return from a factor-free
three/two frame to the preceding exponent-one state. A one-shot initialization
then allows the typed normalizer to be entered without treating a larger
temporary inner source as a source decrease. The resulting well-founded marked
normalization excludes all draw positions.
\end{abstract}

\tableofcontents

\section{The game and the theorem}

The current state is a positive odd integer $n$. For $n>1$ the player to move
chooses one of
\[
  3n-1,\qquad 3n+1,
\]
and divides the chosen value by the largest possible power of $2$, leaving an
odd integer. If the resulting odd integer is $1$, the player who made the move
wins immediately. Players alternate with perfect information.

A position is called \WIN\ for the player to move if at least one move leads to
\LOSS. It is called \LOSS\ if every legal move leads to \WIN. These two classes
are the inductively generated finite-remoteness classes. Any position not in
either class is called \DRAW.

\begin{theorem}[Main theorem]\label{thm:main}
Every positive odd starting position has finite remoteness. Equivalently, the
$3n\pm1$ game has no \DRAW\ positions, and optimal play reaches $1$ after
finitely many moves.
\end{theorem}

The theorem is not a statement about arbitrary play. There are legal cycles,
such as
\[
  5\longrightarrow 7\longrightarrow 5,
\]
under suitable choices by the players.

\section{Binary conjugation}

Write an odd state as $n=2m+1$ and define
\[
  F(m)=\left\lceil\frac{3m}{2}\right\rceil .
\]
For a nonnegative integer $x$, let $R(x)$ be the integer obtained by deleting
the maximal alternating suffix of the binary expansion of $x$. Thus
$R(1001_2)=10_2$ and $R(10101_2)=0$.

\begin{proposition}[First normal form]\label{prop:first-normal}
For $m>0$, the two children in $m$-coordinates are exactly
\[
  F(m),\qquad F(R(m)).
\]
\end{proposition}

\begin{proof}
One has
\[
  3(2m+1)-1=2(3m+1),\qquad
  3(2m+1)+1=2(3m+2).
\]
Exactly one of $3m+1$ and $3m+2$ is odd. In the other expression, each
additional factor of two removed corresponds to crossing one additional
alternating bit at the end of the binary expansion of $m$. The process stops
at the first repeated adjacent bit, which is precisely the deletion defining
$R(m)$. Returning to $m$-coordinates gives the displayed pair.
\end{proof}

Since both children lie in the image of $F$, we represent the state $F(q)$ by
$q$. The conjugated game has terminal state $0$ and moves
\[
  A(q)=\left\lceil\frac{3q}{2}\right\rceil,
  \qquad
  B(q)=R(A(q)).
\]

\begin{proposition}[Strict contraction]\label{prop:B-contracts}
For every $q>0$,
\[
  B(q)<q.
\]
\end{proposition}

\begin{proof}
Deleting a nonempty binary suffix gives
\[
  B(q)\le \left\lfloor\frac{A(q)}2\right\rfloor
       \le \left\lfloor\frac{3q+1}{4}\right\rfloor<q
\]
for $q\ge2$, while $B(1)=0$.
\end{proof}

Thus all unbounded difficulty comes from the expanding branch $A$ and from the
unbounded alternating suffix removed by $B$.

\section{Constant-tail coordinates and coefficient sources}

For a positive odd integer $a$, an exponent $r\ge1$, and a phase
$e\in\{0,1\}$, put
\[
  Q_r^e(a)=a2^r-e.
\]
This is a binary word whose final constant block has length $r$ and consists
of the bit $e$.

\begin{lemma}[Long-tail recurrence]\label{lem:long-tail}
For every positive odd $a$, phase $e$, and $r\ge3$,
\[
  A(Q_r^e(a))=Q_{r-1}^e(3a),\qquad
  B(Q_r^e(a))=Q_{r-2}^e(3a).
\]
Moreover
\[
  B(Q_1^e(a))=B(Q_2^e(a)).
\]
\end{lemma}

\begin{proof}
For $r\ge3$, the expanding child still ends in at least two equal bits; hence
its maximal alternating suffix consists only of the final bit. This gives both
long-tail identities. At the boundary, the binary words of $3a-e$ and $6a-e$
differ by appending one bit that extends the same alternating suffix, so the
same prefix remains after deletion.
\end{proof}

Define the embedded three-free coefficient
\[
  J(s)=2A(s)+1.
\]
Every positive odd coefficient $a$ has a unique factorization
\[
  a=3^kJ(s),\qquad k\ge0.
\]
The integer $s$ is called the coefficient source of $a$ and is denoted
$\src(a)$.

\begin{lemma}[Source boundary]\label{lem:source-boundary}
For every $s>0$ there is a phase
\[
  \alpha(s)=1-\left(\left\lfloor\frac{s}{2}\right\rfloor\bmod2\right)
\]
such that
\[
  \odd(3J(s)+1-2\alpha(s))=J(A(s))
\]
with $2$-adic valuation one, while the opposite phase selects $J(B(s))$ with
valuation at least two.
\end{lemma}

Multiplying a constant-tail coefficient by $3$ preserves its coefficient
source. Passing through a factor-free signed boundary with coefficient $J(s)$
therefore exposes exactly one ordinary $A/B$ move on the source.

\section{Finite proof tokens}

The terminal state $0$ is \LOSS\ with height $h(0)=0$. A finite \WIN\ state has
height
\[
  h(x)=1+\min\{h(y):y\text{ is a \LOSS\ child of }x\},
\]
and a finite \LOSS\ state has height
\[
  h(x)=1+\max\{h(y):y\text{ is a child of }x\}.
\]
Only finite states whose outcome and height comparison have actually been
established are inserted into the marked configuration. They are called proof
tokens.

\begin{lemma}[Multiset token rank]\label{lem:token-rank}
For a finite multiset $M$ of proof heights define
\[
  \Phi(M)=\mathop{\#}_{h\in M}\omega^h,
\]
the natural ordinal sum of the monomials $\omega^h$. Replacing one token of
height $H$ by any finite multiset of certified descendants of heights strictly
below $H$ strictly decreases $\Phi$.
\end{lemma}

\begin{proof}
In Cantor normal form the coefficient of $\omega^H$ decreases by one, and all
inserted terms have lower exponent. No finite collection of lower terms can
compensate for the loss at exponent $H$.
\end{proof}

\section{Factor removal at fixed tail exponent}

The next lemma is the entry device that prevents a factorful exponent-one
state from creating an unranked restart.

\begin{lemma}[Fixed-tail factor removal]\label{lem:fixed-tail-factor}
Let $c$ be positive and odd, $e\in\{0,1\}$, and $r\ge2$. Put
\[
  X=Q_r^e(3c),\qquad G=Q_r^e(c),\qquad H=Q_{r+1}^e(c).
\]
If $X$ is \DRAW, then either $G$ is \DRAW, or the local configuration supplies
finite outcome provenance suitable for a marked typed entry. More explicitly,
if $Y=A(G)$, then
\[
  \Moves(H)=\{X,Y\},
\]
and, whenever $G$ is not \DRAW, at least one finite \WIN/\LOSS\ token is forced
among $G,Y$ and the other child of $G$.
\end{lemma}

\begin{proof}
Since $r+1\ge3$, Lemma~\ref{lem:long-tail} gives
\[
  A(H)=Q_r^e(3c)=X,
  \qquad
  B(H)=Q_{r-1}^e(3c).
\]
For $r\ge3$ this second state is $A(G)$ by the same long-tail recurrence. For
$r=2$, $G=Q_2^e(c)$ and $A(G)=Q_1^e(3c)$, again equal to $B(H)$. Thus
$\Moves(H)=\{X,Y\}$.

Because $X$ is \DRAW, the state $H$ cannot be \LOSS. If $H$ is \WIN, then its
other child $Y$ is \LOSS, giving a finite token. Assume therefore that $H$ is
\DRAW. Then $Y$ cannot be \LOSS. If $G$ is \DRAW, we are in the first
alternative.

It remains to suppose that $G$ is finite. If $G$ is \LOSS, then all its
children, including $Y$, are \WIN. If $G$ is \WIN, then either $Y$ is a \LOSS\
witness or, since $Y$ is not \LOSS, the other child of $G$ is a \LOSS\ witness.
In all cases the displayed local diamond provides an actual finite outcome
token with proved incidence.
\end{proof}

\begin{corollary}[Removing all factor powers at exponent two]\label{cor:remove-to-base}
Let $a=J(s)$ and $k\ge1$. If
\[
  Q_2^e(3^ka)
\]
is \DRAW, then either a finite provenance token is obtained, or
\[
  Q_2^e(a)
\]
is \DRAW.
\end{corollary}

\begin{proof}
Apply Lemma~\ref{lem:fixed-tail-factor} repeatedly with $r=2$ and
$c=3^{i-1}a$. If the token alternative never occurs, the factor exponent is
reduced by one at each step until it reaches zero.
\end{proof}

\begin{lemma}[Exponent-one lift to exponent two]\label{lem:exp-one-lift}
Let $a$ be positive and odd and $k\ge1$. If
\[
  R_k=Q_1^e(3^ka)
\]
is \DRAW, then either
\[
  P_{k-1}=Q_2^e(3^{k-1}a)
\]
is \DRAW, or $P_{k-1}$ is \WIN\ and its other child is a finite \LOSS\ token.
\end{lemma}

\begin{proof}
The long-tail boundary identities give
\[
  A(P_{k-1})=R_k.
\]
Since $P_{k-1}$ has a \DRAW\ child, it cannot be \LOSS. If it is not \DRAW,
it is \WIN; its \DRAW\ child $R_k$ cannot witness that it is \WIN, so its other
child is \LOSS.
\end{proof}

\begin{proposition}[Arbitrary factorful exponent-one entry]\label{prop:factorful-entry}
Let $a=J(s)$ and $k>0$. If
\[
  Q_1^e(3^ka)
\]
is \DRAW, then after finitely many outcome-compatible local steps either a
finite provenance token has been produced, or
\[
  Q_2^e(a)
\]
is \DRAW.
\end{proposition}

\begin{proof}
Apply Lemma~\ref{lem:exp-one-lift}. If it produces a finite token, stop. If it
produces the exponent-two \DRAW\ state $Q_2^e(3^{k-1}a)$, apply
Corollary~\ref{cor:remove-to-base}. This either produces a finite provenance
token or removes all remaining factors of three at the same exponent two.
\end{proof}

\section{Factor-free exponent-two base entry}

It remains to analyze the base state produced by
Proposition~\ref{prop:factorful-entry}. Let $x>0$, put
\[
  a=J(x),\qquad P=Q_1^e(a),\qquad U=Q_2^e(a),
\]
\[
  b=B(P)=B(U),\qquad F=A(U)=Q_1^e(3a),\qquad g=1-e.
\]

Suppose $U$ is \DRAW\ and $x$ is the globally least coefficient source of a
\DRAW. The common child $b$ has canonical coefficient strictly below $J(x)$.
Since $a=J(x)$ is factor-free, this implies that the coefficient source of $b$
is below $x$. Therefore $b$ is not \DRAW. As a child of a \DRAW\ state it also
cannot be \LOSS; hence
\[
  b\text{ is \WIN},\qquad F\text{ is \DRAW}.
\]
The exact first-factor identity is
\[
  \boxed{9J(x)+1-2e=2^vJ(b)}. \tag{1}
\]
Consequently the signed child of $F$ is
\[
  A(F)=Q_v^g(J(b)).
\]

\begin{proposition}[Base-entry attachment]\label{prop:base-entry}
Every factor-free exponent-two \DRAW\ at a globally minimal coefficient source
has a finite continuation that either strictly lowers the retained source,
provides a finite proof-token descent, or enters the typed obligation
normalizer through the one-shot entry component.
\end{proposition}

\begin{proof}
We split according to the phase at $x$.

First suppose $e=1-\alpha(x)$, the $B$-selecting phase. The first source
boundary numerator is divisible by $4$, and
\[
  9J(x)+1-2e=3(3J(x)+1-2e)-2(1-2e)\equiv2\pmod4.
\]
Thus $v=1$ in (1). Writing $u=A(x)$ and using the parity defining the
$B$-selecting phase gives the exact endpoint
\[
  b=3A(x)+1.
\]
The two children of $F$ are
\[
  T=Q_1^g(J(b)),\qquad C=R(T),
\]
and the signed-prefix identity identifies $C$ as the ordinary child of $b$
selected by phase $g$:
\[
  C\in\{A(b),B(b)\}.
\]
At least one of $T,C$ is \DRAW. If $T$ is \DRAW, the first alternative of the
obligation $O(b,g)$ holds and the typed normalizer is initialized. If $C$ is
\DRAW, the other ordinary child of the \WIN\ state $b$ is a genuine \LOSS\
token. For nonexceptional $b$, the ordinary side diamond gives a lower boundary
\WIN\ token. For the exceptional $A$-child orientation, the exceptional side
formula gives an adjacent \DRAW\ frame over the actual \LOSS\ token. In the
sole exceptional $B$-child orientation, the estimate
\[
  C=B(b)\le \frac{A(b)}8\le \frac{3b+1}{16}
\]
combined with $b\le(9x+5)/2$ yields
\[
  \src(C)<\frac{C}{6}\le\frac{3b+1}{96}
  \le\frac{27x+17}{192}<x,
\]
so the retained source strictly decreases.

Now suppose $e=\alpha(x)$, the $A$-selecting phase. Lemma~\ref{lem:source-boundary}
gives
\[
  3J(x)+1-2e=2J(A(x)),
\]
and (1) therefore has $v\ge2$. If $v\ge4$, then
\[
  16J(b)\le 9J(x)+1\le 27x+19.
\]
Since $J(b)\ge3b+1$, we get
\[
  48b+16\le27x+19,
  \qquad
  b\le\frac{9x+1}{16}<x.
\]
Thus any \DRAW\ child of $F$ lies below the retained source.

If $v=2$, the two children of $F$ are exactly
\[
  Q_2^g(J(b)),\qquad Q_1^g(J(b)),
\]
so the obligation $O(b,g)$ holds. If $v=3$, the children form the frame
\[
  Q_3^g(J(b)),\qquad Q_2^g(J(b)).
\]
Reducing (1) modulo $3$ gives
\[
  J(b)\equiv1+g\pmod3,
\]
which is precisely the hidden-parent congruence required for the typed
three/two frame. That frame is therefore normalized by the same source/token
routing used for typed factor frames.
\end{proof}

\section{Typed fixed-fibre normalization}

The local source, factor, and proof-token diamonds used after a typed entry form
a finite routing system. Its control states include $A$-selecting obligations,
$B$-selecting transfers, finite factor forks, high returns, and marked terminal
rows. The important distinction is between retained data and routing data.

The retained data are:
\[
  (\text{outer source},\; \Phi(M),\;\eta),
\]
where $M$ is the finite multiset of carried proof-token heights and
$\eta\in\{1,0\}$ is ordered by $1>0$. The value $\eta=1$ means that the current
outer source/token fibre has not yet initialized the typed normalizer; the
value $\eta=0$ means that the typed normalizer is active.

Routing data include the current source cursor, phase, factor exponent, tail
exponent, and local control state. These quantities may increase numerically.
They do not replace the retained source unless a cited local lemma proves a
strict source comparison.

\begin{proposition}[Typed normalizer]\label{prop:typed-normalizer}
For every typed entry supplied with a retained arithmetic source anchor and the
finite proof-token provenance required by the local diamonds, the equal-retained
data routing relation is well-founded.
\end{proposition}

\begin{proof}[Proof sketch]
The finite routing system has only the following non-strict transitions inside a
fixed retained fibre: bounded canonical $A$-streaks, valuation-two $B$-transfers
with an installed witness, and factor-tail recurrences that carry an existing
finite token. The canonical $A$-streak exits after at most four side tests, since
four consecutive \WIN\ states cannot occur on one $A$-ray. The valuation-two
recycle replaces its witness by a certified lower descendant whenever it repeats.
Marked factor tails either lower their tail counter or replace an existing
\WIN/\LOSS\ token by lower token(s). Hence there is no equal-rank cycle in a
fixed retained fibre. Strict source or token transitions are handled outside the
fibre by the lexicographic rank.
\end{proof}

\begin{lemma}[One-shot entry]\label{lem:eta}
A transition that changes $\eta$ from $1$ to $0$ is strict regardless of the
numerical value of the newly initialized inner routing source. Inside an
unchanged outer source/token fibre, $\eta$ never changes from $0$ back to $1$.
A later raw reset is allowed only after a strict earlier source decrease or a
strict proof-token replacement.
\end{lemma}

\begin{proof}
Use the lexicographic order
\[
  (\text{outer source},\Phi(M),\eta,\text{inner typed rank}).
\]
The component $\eta:1\to0$ strictly decreases after the retained source and
outer token components have been fixed. Because $\eta$ is placed before all
inner routing data, the new inner source may be larger than the old outer source
without being counted as a source update. A change $0\to1$ in the same outer
fibre would increase the rank, so it is not a permitted routing transition.
Only a strict decrease in an earlier component may reset later components.
\end{proof}

\section{Proof of the main theorem}

Assume for contradiction that a \DRAW\ exists and choose the least coefficient
source $s$ of a \DRAW. The source-zero case is handled by the explicit
zero-source constant-tail normalization: a zero-source \DRAW\ reaches an
exponent-one/two boundary with either a strict source exit or a typed finite
terminal token.

Let $s>0$. Take an actual \DRAW\ whose coefficient source is $s$. Long-tail
recurrence reduces it to tail exponent one or two after finitely many
outcome-compatible moves. If a factorful exponent-one state
\[
  Q_1^e(3^kJ(s)),\qquad k>0,
\]
is reached, Proposition~\ref{prop:factorful-entry} gives either finite
provenance for a one-shot typed entry or the factor-free exponent-two state
\[
  Q_2^e(J(s))\text{ \DRAW}.
\]
The latter is handled by Proposition~\ref{prop:base-entry}.

Thus every marked macrostep from a \DRAW\ configuration has one of the following
forms:
\begin{enumerate}[label=(\roman*)]
\item the retained coefficient source strictly decreases;
\item at fixed retained source, an existing finite proof token is replaced by
      certified lower token(s), so Lemma~\ref{lem:token-rank} strictly lowers
      $\Phi$;
\item the retained source and token multiset are unchanged, and the one-shot
      component changes $\eta:1\to0$;
\item $\eta=0$ and the path remains inside the well-founded typed fibre of
      Proposition~\ref{prop:typed-normalizer}.
\end{enumerate}
The lexicographic rank
\[
  (\text{retained source},\Phi(M),\eta,\text{typed fibre rank})
\]
is therefore well-founded.

On the other hand, a \DRAW\ state has no \LOSS\ child and at least one \DRAW\
child. Hence after each finite macrostep one may choose an outcome-compatible
continuing \DRAW\ unless a strict smaller-source contradiction has already
occurred. Repeating this choice would produce an infinite descending path in the
well-founded marked relation, impossible. Hence the conjugated game has no
\DRAW\ positions.

The conjugacy of Section~2 sends the terminal conjugated state $0$ to the
original odd state $1$. Therefore every positive odd starting position has finite
remoteness and optimal play reaches $1$.

\section{Verification boundary and reproducibility}

The proof above is a mathematical proof. The accompanying computational files
are supporting checks with a narrower role:
\begin{itemize}[leftmargin=*]
\item finite Python regressions check the arithmetic identities, valuation rows,
      source inequalities, and finite outcome certificates on stated finite
      ranges;
\item the symbolic global-routing certificate checks the declared typed
      control/rank assembly, but it is not an end-to-end formal proof of the
      theorem;
\item the Lean files in the repository kernel-check selected definitions and
      general well-founded metatheory, not the complete concrete no-\DRAW\
      theorem.
\end{itemize}
A reproducible check should record the repository revision, run the standard
Python check script, run the accompanying arithmetic regression, and read the formal
coverage file before interpreting the machine outputs.

\begin{thebibliography}{9}

\bibitem{Guy1986}
R. K. Guy,
John Isbell's Game of Beanstalk and John Conway's Game of Beans-Don't-Talk,
\emph{Mathematics Magazine} 59(5) (1986), 259--269.

\bibitem{Guy1996}
R. K. Guy,
Unsolved Problems in Combinatorial Games, Problem 42,
in \emph{Games of No Chance}, MSRI Publications 29, Cambridge University Press,
1996.

\bibitem{Althofer2024}
I. Alth\"ofer, M. Hartisch, T. Zipproth,
Analysis of a Collatz Game and Other Variants of the $3n+1$ Problem,
\emph{Advances in Computer Games} (2024), 123--132.

\bibitem{AlthoferPage}
I. Alth\"ofer,
Collatz Problems: The Two-Player $3n+-1$ Game,
\url{https://althofer.de/collatz-prizes.html}.

\bibitem{Repo}
G. Pochuev,
Optimal $3n\pm1$ Game: Proof and Verification Repository,
\url{https://github.com/Grisha-Pochuev/3n-plus-minus-1-game}.

\end{thebibliography}

\end{document}
